%%%%%%%%%%%%%%%%%%%%%%%%%%%%%%%%%%%%%%%%%%%%%%%%%%%%%%%%%%
\begin{frame}
  \begin{center}
    {\Large Understanding Your Data}
  \end{center}
\end{frame}

%%%%%%%%%%%%%%%%%%%%%%%%%%%%%%%%%%%%%%%%%%%%%%%%%%%%%%%%%%
\begin{frame}[fragile]\frametitle{Why Look at Data Before Modeling?}
\begin{itemize}
\item Every ML pipeline so far in this course has started with a dataset -- but real datasets
  arrive messy: missing values, wrong types, outliers, duplicate rows.
\item A model trained on unexamined data doesn't fail loudly -- it just quietly learns the
  wrong thing, or learns noise instead of signal.
\item ``Garbage in, garbage out'' -- no algorithm, however sophisticated, can fix what bad
  input data already broke.
\end{itemize}
\begin{block}{Intuition}
Would you start machining a part from a raw casting without first inspecting it for cracks,
checking its dimensions, and cleaning off the mold sand? Data is the raw material of ML --
skipping inspection just moves the failure downstream, where it's far more expensive to find.
\end{block}
\end{frame}

%%%%%%%%%%%%%%%%%%%%%%%%%%%%%%%%%%%%%%%%%%%%%%%%%%%%%%%%%%
\begin{frame}[fragile]\frametitle{What is Exploratory Data Analysis (EDA)?}
\begin{itemize}
\item EDA is the practice of looking at a dataset -- through summaries, tables, and plots --
  \emph{before} building any model, to understand what you actually have.
\item It answers questions like: What does each column mean? What range of values is normal?
  Are any values missing or clearly wrong? Which variables seem related to each other?
\item EDA is detective work, not decoration -- every chart or table should be answering a
  specific question about the data, not just filling a slide.
\end{itemize}
\end{frame}

%%%%%%%%%%%%%%%%%%%%%%%%%%%%%%%%%%%%%%%%%%%%%%%%%%%%%%%%%%
\begin{frame}[fragile]\frametitle{Two Angles: Univariate and Bivariate}
\begin{itemize}
\item \textbf{Univariate analysis} looks at one column at a time: What's its distribution?
  Its mean, median, spread? Is it skewed? Does it have outliers?
\item \textbf{Bivariate / multivariate analysis} looks at relationships \emph{between}
  columns: Do two numeric columns move together (correlation)? Does a categorical column
  split a numeric one into different groups?
\end{itemize}
\begin{block}{Intuition}
Univariate analysis is like inspecting one machine part on its own -- checking its tolerances
in isolation. Bivariate analysis is checking how two parts fit together -- their relationship
is where a lot of real engineering problems (and real ML signal) actually live.
\end{block}
\end{frame}

%%%%%%%%%%%%%%%%%%%%%%%%%%%%%%%%%%%%%%%%%%%%%%%%%%%%%%%%%%
\begin{frame}[fragile]\frametitle{What EDA Is Looking For}
\begin{itemize}
\item \textbf{Missing data} -- which columns have gaps, and are they random or systematic?
\item \textbf{Outliers} -- values far outside the normal range; sometimes real, sometimes
  sensor glitches or data-entry errors.
\item \textbf{Wrong types} -- a numeric column read in as text, a date stored as a string.
\item \textbf{Imbalance} -- for a target column, are the classes roughly even, or is one
  outcome rare (as with most failure/fraud/churn-style targets)?
\end{itemize}
\end{frame}

%%%%%%%%%%%%%%%%%%%%%%%%%%%%%%%%%%%%%%%%%%%%%%%%%%%%%%%%%%
\begin{frame}[fragile]\frametitle{Modern EDA: Automated Profiling}
\begin{itemize}
\item Everything on the previous slide -- missing data, outliers, types, imbalance -- can be
  generated automatically as a first-pass report, before writing any manual pandas code.
\item Tools like \texttt{ydata-profiling} (formerly pandas-profiling) and \texttt{sweetviz}
  scan a DataFrame and produce a full HTML report: per-column distributions, correlations,
  missing-value maps, warnings about high cardinality or skew.
\item Useful as a fast starting point, not a replacement for judgment -- it flags what to look
  at, it doesn't decide what matters for your specific problem.
\end{itemize}
\begin{lstlisting}[language=Python]
from ydata_profiling import ProfileReport

profile = ProfileReport(df, title="Churn Dataset Report")
profile.to_file("report.html")
\end{lstlisting}
\end{frame}

%%%%%%%%%%%%%%%%%%%%%%%%%%%%%%%%%%%%%%%%%%%%%%%%%%%%%%%%%%
\begin{frame}[fragile]\frametitle{From Exploring to Preparing}
\begin{itemize}
\item Once EDA surfaces a problem -- missing values, a skewed feature, an unencoded category
  -- \textbf{data preparation} is the set of techniques used to fix it before modeling.
\item The two are a loop, not a strict sequence: you explore, prepare, then often explore
  again to confirm the fix worked.
\item The rest of this session works through that toolkit -- handling missing values,
  scaling, encoding, and engineering new features -- then closes with one dataset walked
  through start to finish.
\end{itemize}
\end{frame}
